\documentclass[10pt,twocolumn]{article}
\usepackage[T1]{fontenc}
\usepackage[utf8]{inputenc}
\usepackage{lmodern}
\usepackage[margin=1.8cm,top=2.4cm,bottom=2cm,columnsep=0.6cm]{geometry}
\usepackage{fancyhdr}
\usepackage{titlesec}
\usepackage{booktabs}
\usepackage{amsmath,amssymb,amsfonts}
\usepackage{microtype}
\usepackage{xcolor}
\usepackage{mdframed}
\usepackage{parskip}
\usepackage{enumitem}
\usepackage{hyperref}

\definecolor{rulered}{RGB}{180,30,30}
\definecolor{boxgray}{RGB}{245,245,245}
\definecolor{keywordgray}{RGB}{100,100,100}

\pagestyle{fancy}
\fancyhf{}
\fancyhead[L]{\textsc{\small Claude Mag}}
\fancyhead[C]{\small\textit{Machine Learning Research Digest}}
\fancyhead[R]{\small 2026-07-01}
\fancyfoot[C]{\small\thepage}
\renewcommand{\headrulewidth}{0.8pt}

\titleformat{\section}{\normalsize\bfseries\color{rulered}}{}{0pt}{}%
  [\vspace{-4pt}\color{rulered}\rule{\columnwidth}{0.4pt}]
\titlespacing{\section}{0pt}{8pt}{4pt}

\hypersetup{colorlinks=true,urlcolor=rulered,linkcolor=black}

\begin{document}
\setlength{\parindent}{0pt}
\setlength{\parskip}{4pt}

{\small\color{keywordgray}\textbf{Keywords:}
  offline reinforcement learning \textbullet{}
  flow matching \textbullet{}
  Q-learning \textbullet{}
  diffusion policy}\par\vspace{4pt}

{\LARGE\bfseries\raggedright
  Reversal Q-Learning}\par\vspace{4pt}

{\large\itshape\raggedright
  Flow-Based Offline RL Trains Expressive Policies Without Backpropagating Through
  the Denoising Chain}\par\vspace{6pt}

{\small\textbf{By} Aditya Oberai, Seohong Park \& Sergey Levine
  (UC Berkeley)}\par\vspace{2pt}

{\small\color{keywordgray}arXiv:2606.17551 \textbullet{} Submitted June 2026}
\par\vspace{2pt}\noindent{\color{rulered}\rule{\columnwidth}{0.6pt}}\vspace{6pt}

Flow-matching policies can now be trained from offline data without differentiating
through their iterative denoising chains. \textit{Reversal Q-learning} (RQL) treats
each denoising step as an MDP action, reverses completed flows to generate synthetic
on-policy trajectories, and applies standard Q-learning---beating prior flow-based
offline RL methods on every D4RL locomotion task while eliminating the instability
of backpropagation through time.

\begin{mdframed}[backgroundcolor=boxgray,linecolor=rulered,linewidth=1pt,
    innerleftmargin=6pt,innerrightmargin=6pt,
    innertopmargin=6pt,innerbottommargin=6pt]
\textbf{\small AT A GLANCE}\par\vspace{4pt}
\begin{description}[leftmargin=0pt,labelsep=4pt,itemsep=2pt,parsep=0pt,
    font=\small\bfseries]
  \item[Problem:] Training an expressive flow-matching policy with offline RL
    requires either weak guidance distillation or expensive, unstable backpropagation
    through the full denoising chain.
  \item[Why it matters:] Multimodal action distributions are common in robotics and
    navigation datasets; flow policies can represent them but existing RL training
    methods are ill-suited to offline settings.
  \item[Method:] Treat denoising steps as MDP actions, reverse completed flows to
    obtain on-policy trajectories from offline data, and train with standard
    off-policy Q-learning plus a bias-variance reduction step.
  \item[Result:] State-of-the-art on 12 D4RL locomotion tasks with faster and
    stabler training than BPTT-based methods.
  \item[Evidence:] Systematic evaluation on D4RL medium, medium-replay,
    and medium-expert splits vs.\ IQL, TD3+BC, DQL, IDQL, SfBC.
\end{description}
\end{mdframed}

\section{The Big Picture}

Offline reinforcement learning trains a policy from a fixed dataset of past
experience without additional environment interaction. The challenge is
simultaneously estimating a value function over out-of-distribution actions
(hard) and representing multi-modal action distributions that arise when the
dataset mixes several different behavioral strategies (also hard).

Flow matching---a continuous normalising flow that transforms noise into an
action via an ordinary differential equation (ODE)---is a natural policy class
for the multi-modal problem. Its generation process is smooth, invertible, and
arbitrarily expressive. But plugging flow policies into offline RL has been
awkward: distillation-based methods (SfBC, IDQL) under-use the value function,
and direct RL fine-tuning via policy gradient requires differentiating through
all $T$ ODE steps, which is computationally expensive and numerically fragile.

RQL resolves this by lifting the problem into an \emph{expanded} Markov decision
process where each denoising step is its own atomic action. Standard
off-policy Q-learning then suffices---no differentiation through the ODE solver,
full use of the Q-function at every step.

\section{Why Existing Approaches Fall Short}

\textbf{Guidance distillation} (e.g., IDQL, SfBC) trains a diffusion model via
behaviour cloning and then biases sampling towards high-Q regions at test time.
The policy is never directly trained to maximise $Q$; it remains anchored to the
behavioural prior.

\textbf{BPTT fine-tuning} (e.g., DQL) back-propagates a policy gradient through
the full $T$-step ODE unrolling. This suffers from (i) exploding/vanishing
gradients over long chains, (ii) prohibitive memory cost, and (iii) biased
gradients because the ODE solver itself is not differentiable in the strict sense.

Neither exploits the structure of the flow ODE for off-policy learning.

\section{Core Mechanism}

\textbf{Expanded MDP.}
Discretise $[0,1]$ into $T$ steps with $\Delta t=1/T$. Define an augmented
state $\tilde s_k=(s,x_{k\Delta t})$ where $s$ is the environment state and
$x_k$ is the current latent. The action at step $k$ is $\delta_k =
v_\theta(\tilde s_k)\cdot\Delta t$. Reward is $R_T = Q^*(s,x_1)$ at the
final step and zero elsewhere.

\textbf{Reversal.}
Given an offline pair $(s,a^*)\in\mathcal D$, reconstruct the full latent
trajectory by running the ODE \emph{backwards}:
$x_{k-1}^* = x_k^* - v_\theta(x_k^*,k\Delta t; s)\cdot\Delta t$,
with $x_T^* = a^*$. This converts each offline sample into a complete
expanded-MDP episode.

\textbf{Q-learning.}
Apply Bellman updates over these synthetic trajectories. The flow policy is
updated by maximising $Q_\phi(\tilde s_k, v_\theta(\tilde s_k))$ at each
step---a single-step gradient, never multi-step.

\section{Technical Formulation}

A flow matching policy defines an ODE:
\[
\frac{dx_t}{dt} = v_\theta(x_t, t; s), \quad x_0 \sim \mathcal{N}(0,I),
\quad a = x_1.
\]

The expanded-MDP Q-function satisfies:
\[
Q^*(\tilde s_k, \delta_k) =
\begin{cases}
  Q^*(s, x_{k+1}) & k < T-1\\[2pt]
  r + \gamma\,\mathbb{E}_{s'}[V^*(s')] & k = T-1
\end{cases}
\]

The Q-network is trained with:
\[
\mathcal{L}_Q(\phi) =
\mathbb{E}\!\left[\bigl(Q_\phi(\tilde s_k,\delta_k^*) - y_k\bigr)^2\right]
\]
and the policy is updated by:
\[
\mathcal{L}_\pi(\theta) =
-\mathbb{E}_{k,s,x_0}\!\left[Q_\phi\!\left(\tilde s_k,\,
v_\theta(\tilde s_k)\right)\right].
\]

Crucially, gradients of $\mathcal{L}_\pi$ pass only through the
\emph{single-step} output $v_\theta(\tilde s_k)$, not through the ODE solver.

\section{Learning or Inference Procedure}

\begin{enumerate}[itemsep=2pt,parsep=0pt]
  \item Sample $(s,a^*)\in\mathcal D$.
  \item Run the reverse ODE to obtain $\{x_k^*\}_{k=0}^T$.
  \item Compute Bellman targets $\{y_k\}$ using current $Q_\phi$.
  \item Update $Q_\phi$ via $\mathcal{L}_Q$.
  \item Update $v_\theta$ via $\mathcal{L}_\pi$ (single-step gradient only).
  \item At test time: run the forward ODE from $x_0\sim\mathcal{N}(0,I)$
    for $T$ steps to generate action $a=x_1$.
\end{enumerate}

\section{What the Guarantee Says}

The paper proves that, under the expanded MDP, the optimal Q-function $Q^*$
decomposes across denoising steps exactly as stated, and that the reversal
procedure yields unbiased estimates of the on-policy return. A bias-variance
reduction technique---weighting intermediate step returns by the learned
Q-function rather than bootstrapping all the way from the final step---reduces
variance by a factor proportional to $T$ without introducing bias.

\section{Experimental Findings}

Evaluated on the 12 D4RL locomotion tasks (HalfCheetah, Hopper, Walker2d at
medium / medium-replay / medium-expert quality), RQL achieves:
\begin{itemize}[itemsep=2pt,parsep=0pt]
  \item Highest average normalised score vs.\ IQL, TD3+BC, DQL, IDQL, SfBC.
  \item Particularly large gains on medium-expert splits, where multi-modal
    data benefits most from an expressive policy.
  \item Stable training curves without the loss spikes observed in DQL (BPTT).
  \item Wall-clock training time comparable to IQL despite operating over an
    expanded $T$-step MDP.
\end{itemize}

\section{Ablations and Interpretation}

\textbf{Number of steps $T$:} Performance plateaus around $T=10$; $T=1$
reduces to standard Q-learning with a one-step denoiser, which underperforms.

\textbf{Reversal vs.\ forward noise:} Replacing the reversal with forward
noising (add Gaussian noise to $a^*$) degrades performance by $\approx8$
normalised score on average, confirming that faithfully reconstructing the
latent trajectory is important.

\textbf{Q-weighting at intermediate steps:} Removing the bias-variance
reduction and using terminal Q-values only increases variance and slows
convergence by $\approx2\times$.

\section*{Reference}

Aditya Oberai, Seohong Park, Sergey Levine. \textbf{Reversal Q-Learning}.
arXiv:2606.17551, June 2026.
\url{https://arxiv.org/abs/2606.17551}

\end{document}
